% ============================================================================
% SEC05.TEX - Section 10.5: Win/Lose Conditions
% ============================================================================

\section{Win/Lose Conditions}

\begin{learningoutcomes}
\item Mengimplementasikan win conditions
\item Mengimplementasikan lose conditions
\item Membuat game over logic
\end{learningoutcomes}

\subsection{Win Condition}

\begin{lstlisting}[language={[Sharp]C}]
public class WinCondition : MonoBehaviour
{
    public int targetScore = 100;
    public GameObject victoryUI;
    
    void Update()
    {
        if (ScoreManager.Instance.GetScore() >= targetScore)
        {
            WinGame();
        }
    }
    
    void WinGame()
    {
        victoryUI.SetActive(true);
        GameManager.Instance.ChangeState(GameState.Victory);
        Time.timeScale = 0f;
    }
}
\end{lstlisting}

\subsection{Lose Condition}

\begin{lstlisting}[language={[Sharp]C}]
public class PlayerHealth : MonoBehaviour
{
    public int maxHealth = 100;
    private int currentHealth;
    public GameObject gameOverUI;
    
    void Start()
    {
        currentHealth = maxHealth;
    }
    
    public void TakeDamage(int damage)
    {
        currentHealth -= damage;
        
        if (currentHealth <= 0)
        {
            Die();
        }
    }
    
    void Die()
    {
        gameOverUI.SetActive(true);
        GameManager.Instance.ChangeState(GameState.GameOver);
        Time.timeScale = 0f;
    }
}
\end{lstlisting}

\begin{catatan}
Win/lose conditions memberikan tujuan dan konsekuensi dalam game. Pastikan UI dan feedback yang jelas untuk player.
\end{catatan}

\subsection{Ringkasan Section}

\begin{itemize}
\item Win condition: Tujuan yang harus dicapai
\item Lose condition: Kondisi yang mengakhiri game
\item Game over logic dengan UI feedback
\item Time.timeScale untuk pause saat game over
\end{itemize}
